\section{Introduction}

% ---------------------------------------------------------------
% OUTLINE
%   P1  clinical motivation (PSMA PET/CT, longitudinal imaging)
%   P2  why deformable registration is hard here + challenge task
%   P3  related work: classical -> learning-based -> LapIRN; rigidity
%   P4  contributions
% >>> FIGURE (teaser, OPTIONAL -- cut first if over length):
%     a baseline/follow-up pair illustrating the FOV, positioning and
%     bladder-filling differences described in P2. If cut, fold the
%     example into the qualitative figure in the Results.
% ---------------------------------------------------------------

Prostate-specific membrane antigen (PSMA) PET/CT has become a central imaging modality for staging, treatment planning, radiopharmaceutical therapy, and post-therapy response assessment in prostate cancer~\cite{ref_psma_clinical}.
In these settings, patients often undergo multiple scans over time, comprising a baseline scan before therapy and one or more follow-up scans afterwards; accurate baseline-to-follow-up registration can support comparison of disease burden, localisation of uptake changes, assessment of lesion progression or response, dose accumulation, and quantitative longitudinal analysis~\cite{ref_psmareg_challenge}.
The challenge organisers identify several factors that make establishing this correspondence by deformable image registration difficult in routine whole-body data~\cite{ref_psmareg_challenge}: baseline and follow-up scans differ substantially in field of view, patient positioning, respiratory state, bladder and bowel filling, and body weight; CT and PET encode complementary but appearance-dissimilar information; and the registration must preserve the quantitative PET signal used for response assessment while remaining physically plausible structures such as the skeleton.
The Learn2Reg~2026 PSMAReg challenge~\cite{ref_learn2reg_alessa,ref_learn2reg_hansen,ref_psmaregdata} formalises this problem as the estimation of a dense displacement field that warps each follow-up (moving) scan to its baseline (fixed) scan, evaluated jointly on organ overlap, boundary distance, deformation regularity, and preservation of PET-derived tumour biomarkers.

Deformable registration has traditionally been posed as an iterative optimisation over a regularised similarity objective, with widely used instances including the diffeomorphic Demons~\cite{ref_demons}, SyN~\cite{ref_syn}, and B-spline~\cite{bspline} schemes.
While accurate, these methods are slow at whole-body resolution.
Learning-based methods instead amortise the optimisation over a training set into the network weights, replacing per-pair iterative optimisation with a single forward pass at inference.
VoxelMorph established the unsupervised paradigm~\cite{ref_voxelmorph_cvpr}, and subsequent work introduced diffeomorphic parameterisations through stationary velocity fields~\cite{ref_svf_voxelmorph}, transformer backbones such as TransMorph~\cite{ref_transmorph}, and multi-resolution architectures such as LapIRN~\cite{ref_lapirn}.
Among the latter, LapIRN uses a Laplacian pyramid of coarse-to-fine sub-networks to recover large deformations while keeping the transformation smooth and invertible, and has been a strong performer across previous Learn2Reg editions~\cite{ref_learn2reg_alessa,ref_learn2reg_hansen}.
A recurring limitation of purely intensity-driven losses is that they permit non-physical folding.
Anatomy-aware regularisation that penalises non-rigid deformation, e.g. of skeletal structures~\cite{ref_rigidity_1_loeckx,ref_rigidity_2_ruan,ref_rigidity_3_staring,ref_rigidity_4_staring_modersitzki}, is therefore particularly relevant for whole-body PET/CT, where the skeleton is a frequent site of PSMA-avid disease.

In this work, we adapt LapIRN to longitudinal whole-body PSMA PET/CT for the PSMAReg task.
We retain its diffeomorphic, coarse-to-fine formulation to handle the large field-of-view shifts between visits, and augment it with several task-specific regularisation terms targeted at the challenge objectives: a bone-rigidity loss that discourages implausible deformation of the skeleton, a total-lesion-glycolysis (TLG) loss that preserves PET-derived lesion burden under the estimated deformation, and a term based on the Jacobian determinant within the tumour volume that acts as a proxy for preserving the metabolic tumour volume (MTV).
We describe the method and the training and inference infrastructure, and report validation leaderboard results for each metric.

% --- Contribution list. PLACEHOLDER: wording of (v) depends on the
% --- final model choice; keep the list to four or five items.
In summary, our contributions are:
\begin{itemize}
  \item an adaptation of the diffeomorphic, coarse-to-fine LapIRN
        architecture to longitudinal whole-body PSMA PET/CT, combining
        an affine pre-registration with a three-level pyramid and
        test-time instance optimisation;
  \item a set of biomarker-preserving losses that constrain the
        PET-derived MTV and TLG under the estimated deformation, including
        a Jacobian-based surrogate that supplies dense gradients
        throughout the lesion interior;
  \item a per-structure rigid-fit (Kabsch) bone-rigidity penalty that,
        unlike derivative-based rigidity terms, reads no voxels outside
        the bone mask and is therefore well posed on thin cortical
        structures;
  \item a gradient-level analysis of the objective showing that the
        apparent accuracy-versus-preservation trade-off is in fact a
        regulariser-versus-data-term tension;
  \item a score-driven checkpoint and submission selection procedure
        that mirrors the challenge ranking. \todo{finalise once the
        model is fixed}
\end{itemize}
